% s-098-p098-m001; original PDF 98, 99
\Aphorism{s-098-p098-m001}{99}
Gedanken mit Gefühlsfarbe bedeuten, dass das Gehirn als Siebvorrichtung nicht engmaschig genug ist, den Antheil niederer Nervenressorts auszusondern. Die reinste Gedankenrasse entspringt aus cerebraler Inzucht, wo Anschauungen, Bilder, endlich Gedanken selbst sich gegenseitig befruchten. „Gefühl“ ist noch ungefilterte Subjectivität, ist die gröbere, auch geschwindere Antwort des Organismus auf Umweltreize: die Süssigkeit der Kirsche ist, subjectiv, eher da als ihre Rundung und Röthe. — Der Mensch hat, selbst nach den Jugendjahren, ganze Epochen, wo er mehr mit den Ganglien denkt als mit dem Gehirn; Gedanken aus solcher Zeit erkennt man daran, dass sie späterhin nicht eigentlich auf eine der beiden Seiten \Sperr{wahr} oder \Sperr{unwahr} zu bringen sind. In ihnen ist viel Energie beisammen, vielleicht die ganze, die der Mensch aufzubieten hatte, aber solche, die im Sinne der Erkenntniss nicht umsetzbar, nicht verwandlungsfähig ist: Wärme, die nicht arbeiten, und Arbeit, die nicht erwärmen will. Der Autor dieses Buches würde sich nicht verzeihen, gegen die eigenen Ganglien zu nachsichtig gewesen zu sein; man wird etwas von einem eisernen Druck spüren, unter dem er sein Allzupersönliches, jugendlich Ekstatisches gehalten hat. Aber die Sclaven wollen ihre Saturnalien; Jugend, Gefühl, Liebe, Metaphysik — die vorintellectuelle Affectenwelt will sich an einer Stelle rednerisch entladen, wenn ihr sonst überall der Mund verboten wird. Seien wir also milde, wie Aristoteles gegen die Tragödie; nehmen wir um der gründlichen Katharsis willen ein wenig Übermass, Dumpfheit, Leidenschaft in den Kauf. Wer ein kühles, sachliches, souverain geistiges Klima liebt, sei vor diesem feuchtschwülen Intermezzo, vor dieser \Foreign{„petite Afrique“} mit ihren Treibhausgewächsen rechtzeitig gewarnt: der Philosoph mache einen Umweg — und der Autor macht ihn gern mit.
\par

